\scanpage{103}
\mspage{47}
% archive sheet number: 102
3) La $W$-équivalence grossière, elle, est ``bonne partout'' -- mais ça dépend pour qui. Pour les \inserted{\deleted{\unsure{tr}} 6} catégories qui sont \deleted{soit à l'intérieur, soit} au dessus de $\underline{\mathrm{Fibhot}}\,S$, i.e.

\[
 \underline{\mathrm{Fibhot}}\,S,\ \underline{\mathrm{Parf}}\,S,\ \underline{\mathrm{Fib}}_0\,S,\ \underline{\mathrm{Fibsc}}_0\,S,\ \underline{\mathrm{Cofib}}_0\,S,\ \underline{\mathrm{Cofibsc}}_0\,S,
\]

elle coïncide avec toutes les autres équivalences \unsure{$\in$|$\varepsilon$} (moyennant l'axiome L5), pour les trois catégories bordantes supérieures gauches

\[
 \underline{\mathrm{Liss}}\,S,\ \underline{\mathrm{Fib}}\,S,\ \underline{\mathrm{Fibsc}}\,S
\]

du diagramme, elle coïncide avec la $W$-équiv. fine droite, et dualement sur

\[
 \underline{\mathrm{Prop}\,S},\ \underline{\mathrm{Cofib}\,S},\ \underline{\mathrm{Cofibsc}}\,S
\]

elle coïncide \deleted{\unsure{avec}} avec la $W$-équivalence fine gauche. Quant à $\mathrm{Cat}_{/S}$ tout entier, c'est la seule des \deleted{quatorze notions} \inserted{catégories du} \deleted{de $W$-équivalence qui y donne une localisé \unsure{distincte}.} diagramme où il y ait \uline{trois} notions intéressantes différentes de $W$-équivalence\unsure{s}, à savoir : la grossière, et les deux fines (à gauche, à droite). Sur les autres \inserted{six} catégories bordantes du diagramme, il y en \deleted{\unsure{a}} en principe deux : la grossière, et une spécifiée des deux fines (défendu de regarder l'autre -- à ses risques et périls \ldots) : la droite du côté gauche, la gauche du côté droit. (NB Je devrais refaire le diagramme, décidément, en mettant les fibrations à droite, ou à gauche.) Quant aux \inserted{six} catégories internes \deleted{\unsure{au}} diagramme, dans \unsure{icelles|elles} il n'y a plus qu'une seule notion de $W$-équivalence (puisque sur elles les quatre notions coïncident).

